\documentclass[a4paper]{article}
\usepackage[utf8]{inputenc}
\usepackage[T1]{fontenc}
\usepackage[margin=1.5cm]{geometry}
\usepackage{tikz}
\usepackage{tikz-uml}
\usepackage{adjustbox}
\tikzumlset{fill class=white, text=black, draw=black}

\begin{document}

\begin{center}
\begin{adjustbox}{max width=\textwidth}
\begin{tikzpicture}

  \umlclass[x=0, y=0]{Utente}{
    - id : int \\
    - nomeUtente : String \\
    - email : String \\
    - password : String
  }{
    + registrati() : void \\
    + login() : boolean \\
    + logout() : void \\
    + scriviRecensione() : void
  }

  \umlclass[x=8, y=0]{Recensione}{
    - id : int \\
    - voto : int \\
    - testo : String \\
    - data : Date
  }{}

  \umlclass[x=16, y=0]{Ricetta}{
    - id : int \\
    - titolo : String \\
    - istruzioni : String \\
    - tempoPreparazione : int \\
    - difficolta : String
  }{
    + visualizzaIstruzioni() : String \\
    + cercaPerIngredienti(lista : Ingrediente[]) : Ricetta[]
  }

  \umlclass[x=16, y=7]{TipoCibo}{
    - nome : String
  }{}

  \umlclass[x=16, y=-6.5]{IngredienteRicetta}{
    - quantita : float \\
    - unita : String
  }{}

  \umlclass[x=8, y=-6.5]{Ingrediente}{
    - id : int \\
    - nome : String
  }{}

  % Relazioni
  \umlassoc[mult1=1, mult2=0..*, arg=scrive]{Utente}{Recensione}
  \umlassoc[mult1=0..*, mult2=1, arg=riguarda]{Recensione}{Ricetta}
  \umlassoc[mult1=0..*, mult2=1..*, arg=è di tipo]{Ricetta}{TipoCibo}
  \umlassoc[mult1=1, mult2=1..*, arg=contiene]{Ricetta}{IngredienteRicetta}
  \umlassoc[mult1=0..*, mult2=1, arg=usa]{IngredienteRicetta}{Ingrediente}

\end{tikzpicture}
\end{adjustbox}
\end{center}

\end{document}